% Reconstructed text from V5 PDF page 21; non-native source
\par\smallskip\noindent\rule{\textwidth}{0.4pt}\par{\sffamily\footnotesize Original V5 page 21 -- recovered text}\par\begin{flushleft}\footnotesize\ttfamily
\par\smallskip
\par\smallskip
\noindent Diagnostics and Serving Infrastructure for MAS 21\par
\par\smallskip
\par\smallskip
\par\smallskip
\noindent management — particularly KV cache eviction policies — can affect safety system effective-\par
\noindent ness when it removes state that safety monitors depend on. We formalize this interaction in\par
\noindent Section 6.\par
\par\smallskip
\noindent 4.4 Benchmarks\par
\noindent Progress in diagnostics depends on benchmarks that vary in failure coverage, trace fidelity, and\par
\noindent evaluation scope. Table 3 summarizes 12 Tier-1 benchmarks.\par
\par\smallskip
\noindent Failure-attribution benchmarks. TraceElephant [6] provides the highest trace fidelity among\par
\noindent surveyed benchmarks: 220 annotated agent failures across three production systems (Captain-\par
\noindent Agent, Magentic-One, SWE-Agent) and three task sources (GAIA, AssistantBench, SWE-\par
\noindent Bench Verified), with 380 total execution traces (§3.1, Table 1). Each trace is a structured JSON\par
\noindent with two layers: trace metadata (task ID, system architecture, agent configuration) and step-\par
\noindent level records (agent ID, input context, output content, tool logs with tool name/inputs/outputs/status)\par
\noindent (§3.3). Annotation was performed by three human annotators with at least one year of MAS de-\par
\noindent bugging experience in a three-round protocol (independent → joint consensus → cross-review),\par
\noindent achieving Krippendorff’s α of 0.72 at the agent level and 0.64 at the step level (§3.4).\par
\noindent  The central quantitative finding is the impact of trace completeness on attribution accuracy.\par
\noindent Under the same All-at-Once attribution protocol, complete traces raise agent-level accuracy\par
\noindent from 51\% (output-only logs matching Who\&When’s observability) to 62.2\% (+11.2 percent-\par
\noindent age points, a 22\% relative gain) (ACL 2026, §4.3, Table 3). At the step level, the gain is more\par
\noindent pronounced: 16\% to 28.1\% (+12.1pp, 76\% relative), confirming that step-level attribution is\par
\noindent substantially more sensitive to missing trace information. Dynamic replay-based attribution\par
\noindent adds a further 10\% step improvement (Table 2). Environment-interacting agents account for\par
\noindent over 50\% of failure attributions, with orchestrators responsible for 18–29\% (§4.2.2, Figure 4).\par
\noindent Average LLM invocations per trace range from 20.5 (Captain-Agent) to 29.3 (Magentic-One)\par
\noindent (§4.2.3). The benchmark is available under CC BY 4.0 with a reproducible execution environ-\par
\noindent ment.\par
\noindent  Who\&When [48] provides 184 failures from 127 multi-agent systems as the first dedicated\par
\noindent failure-attribution benchmark (ICML 2025). The benchmark uses algorithm-generated (via\par
\noindent CaptainAgent) and hand-crafted (via Magentic-One) MAS configurations with up to 5 agents\par
\noindent each, producing log lengths of 5–130 steps (§3.1). Attribution is defined through a “decisive\par
\noindent error” concept (Definitions 2.1, 2.2): the earliest agent action whose correction would change\par
\noindent the task outcome from failure to success. Three annotators with AI-agent expertise performed\par
\noindent an 84.3 human-hour annotation protocol (§3.3), with 15–30\% of instances marked as uncertain.\par
\noindent Best reported accuracy is 53.5\% at the agent level and only 14.2\% at the step level; step-level\par
\noindent accuracy degrades to approximately 0\% on trajectories exceeding 93 steps (§4.4, Figure 4). The\par
\noindent hybrid attribution method costs 149,177 tokens per instance (Table 3), making it one of the most\par
\noindent expensive evaluation protocols.\par
\noindent  Domain-specific alternatives include RootSE (introduced alongside TrajAudit [35], 102\par
\noindent coding failures with step-level annotations from 35 repositories across Python, Go, and JavaScript;\par
\noindent average 51.7 steps per trajectory) and τ -bench (airline and retail customer-service scenarios\par
\noindent used in AgentRx evaluation). The Who\&When Pro benchmark (NeurIPS 2026) extends the\par
\noindent original with 12,326 traces across 26 benchmarks.\par
\par\smallskip
\noindent Reliability and safety benchmarks. ReliabilityBench (1,280 episodes) evaluates pass@1\par
\noindent overestimation by 20–40\% [a], identifying a systematic optimism bias in current evaluation:\par
\end{flushleft}
